% ============================================================================
% 4.4 uart_rx.v
% ============================================================================
\subsection{设计目标}

与发送器相反：持续监视 \texttt{rx\_pin}。当 \texttt{rx\_pin} 从 1 变 0
（起始位到来）时，启动一个计数器。在每个比特的中心位置采样
一次电平，攒够 8 位后输出 \texttt{rx\_data} 并拉高 \texttt{rx\_data\_valid}。

接收器比发送器难，因为：
\begin{itemize}
    \item 发送器是自己决定什么时候发一位，时钟完全可控
    \item 接收器不知道对方何时开头发送，必须自己检测起始位
    \item 接收器不能保证自己的时钟和对方时钟完全对齐，
        必须在比特\textbf{中心}采样以获得最大噪声容限
\end{itemize}

\subsection{状态机设计}

\begin{center}
\begin{tabular}{l l}
\textbf{S\_IDLE} & 监视 rx\_pin。检测到下降沿 → 进入 START 状态。 \\
\textbf{S\_START} & 计数半个周期，再采样一次确认真的是低电平（防毛刺）。 \\
\textbf{S\_REC\_BYTE} & 每过一个周期采样一位，共 8 位，低位在前。 \\
\textbf{S\_STOP} & 等待停止位，完成后回到 IDLE。 \\
\textbf{S\_DATA} & 拉高 valid 标志，让上层模块取走字节。 \\
\end{tabular}
\end{center}

\subsection{源码与关键思路讲解}

\begin{lstlisting}[style=verilog,caption=Chip/src/uart\_rx.v]
module uart_rx
#(
    parameter CLK_FRE = 50,
    parameter BAUD_RATE = 115200
)
(
    input clk,
    input rst_n,
    output reg [7:0] rx_data,       // 收到的字节
    output reg rx_data_valid,       // 一次有效，表明有新字节
    input rx_data_ready,             // 上层已取走，让我进入下一轮
    input rx_pin                     // 串行输入引脚
);

localparam CYCLE = CLK_FRE * 1000000 / BAUD_RATE;
localparam S_IDLE      = 1;
localparam S_START     = 2;
localparam S_REC_BYTE  = 3;
localparam S_STOP      = 4;
localparam S_DATA      = 5;

reg [2:0] state;
reg [2:0] next_state;
reg       rx_d0;           // 打一拍，同步 rx_pin
reg       rx_d1;           // 再打一拍，用于检测下降沿
wire      rx_negedge;      // 检测到一个下降沿
reg [7:0] rx_bits;         // 保存逐位收到的数据
reg [15:0] cycle_cnt;      // 当前位已持续了多久
reg [2:0] bit_cnt;         // 已收到第几位
\end{lstlisting}

\paragraph{打两拍：跨时钟域的基本操作}
\texttt{rx\_pin} 是一个从外部引脚进入的信号，它的变化不一定对齐
我们的时钟沿。直接用它来驱动逻辑会有"亚稳态"风险。
解决办法是\textbf{打两拍}：

\begin{lstlisting}[style=verilog]
assign rx_negedge = rx_d1 && ~rx_d0;  // 检测到从 1→0

always @(posedge clk or negedge rst_n)
begin
    if (!rst_n)
    begin
        rx_d0 <= 1'b0;
        rx_d1 <= 1'b0;
    end
    else
    begin
        rx_d0 <= rx_pin;
        rx_d1 <= rx_d0;
    end
end
\end{lstlisting}

打两拍后，我们内部观察到的信号变化至少延迟了两个时钟周期，
但它是\textbf{稳定的}——不会在一个时钟沿前后出现混乱的值。

\paragraph{状态转移逻辑：使用两个 always 块}
接收器用了"状态 / 次态分离"的写法，这是更标准的状态机写法：

\begin{lstlisting}[style=verilog]
// 第一块：状态更新（时序逻辑，只在 clk 沿更新）
always @(posedge clk or negedge rst_n)
begin
    if (!rst_n)
        state <= S_IDLE;
    else
        state <= next_state;
end

// 第二块：计算 next_state（组合逻辑，异步计算）
always @(*)
begin
    case (state)
        S_IDLE:
            if (rx_negedge)
                next_state = S_START;
            else
                next_state = S_IDLE;
        S_START:
            if (cycle_cnt == CYCLE - 1)
                next_state = S_REC_BYTE;
            else
                next_state = S_START;
        S_REC_BYTE:
            if (cycle_cnt == CYCLE - 1 && bit_cnt == 3'd7)
                next_state = S_STOP;
            else
                next_state = S_REC_BYTE;
        S_STOP:
            if (cycle_cnt == CYCLE/2 - 1)  // 半位就进入数据锁存
                next_state = S_DATA;
            else
                next_state = S_STOP;
        S_DATA:
            if (rx_data_ready)
                next_state = S_IDLE;
            else
                next_state = S_DATA;
        default:
            next_state = S_IDLE;
    endcase
end
\end{lstlisting}

\paragraph{为什么在 S\_STOP 里只等半位}
我们是在起始位的\textbf{中间}检测到稳定低电平的。
后续每位也都是从位中心开始的一个完整周期采样。
但停止位其实不需要等完整的一位——只要确认它是高电平就完成了。
所以用 \texttt{CYCLE/2 - 1}，只等半个周期，
省出的时间让我们更快回到 IDLE 状态，不漏掉下一个字节。

\paragraph{数据与有效标志}

\begin{lstlisting}[style=verilog]
always @(posedge clk or negedge rst_n)
begin
    if (!rst_n)
        rx_data_valid <= 1'b0;
    else if (state == S_STOP && next_state != state)
        rx_data_valid <= 1'b1;
    else if (state == S_DATA && rx_data_ready)
        rx_data_valid <= 1'b0;
end

always @(posedge clk or negedge rst_n)
begin
    if (!rst_n)
        rx_data <= 8'd0;
    else if (state == S_STOP && next_state != state)
        rx_data <= rx_bits;  // 把 8 位并行输出
end
\end{lstlisting}

\paragraph{逐位接收}

\begin{lstlisting}[style=verilog]
always @(posedge clk or negedge rst_n)
begin
    if (!rst_n)
        cycle_cnt <= 16'd0;
    else if ((state == S_START || state == S_REC_BYTE ||
              state == S_STOP) && cycle_cnt < CYCLE - 1)
        cycle_cnt <= cycle_cnt + 16'd1;
    else
        cycle_cnt <= 16'd0;
end

always @(posedge clk or negedge rst_n)
begin
    if (!rst_n)
        rx_bits <= 8'd0;
    else if (state == S_REC_BYTE && cycle_cnt == CYCLE/2 - 1)
        rx_bits[bit_cnt] <= rx_pin;   // 在每个比特的中心采样
end

endmodule
\end{lstlisting}

\paragraph{中心采样的重要性}
在 \texttt{cycle\_cnt == CYCLE/2 - 1} 时采样，相当于在
每个位的正中间读取电平。UART 通信两端的时钟总有微小的误差，
但只要误差不累计到超过半个比特周期，中心采样就可以正确读取数据。
这是 UART 协议的核心鲁棒性来源。

如果在边沿采样，发送器和接收器的时钟漂移累积起来，几个字节之后
你就会开始读错——因为你采样的位置在两个比特的交界处，
读出的电平取决于先到哪一侧。
